\documentclass[12pt]{amsart}
\usepackage{amssymb} % this package is indispensable, among other things, for rendering the \varsubsetneq command
\usepackage{amsmath}
\usepackage{latexsym}
\usepackage{amsfonts}
\usepackage{enumerate}
%\usepackage[mathscr]{eucal}
\usepackage[mathscr]{euscript}
\usepackage{eqlist}
\usepackage{amsthm}
\usepackage{mathtools}
\usepackage{enumitem} % this allows us to create enumerate environments with different item numbering schemes
\usepackage{hyperref} % this allows us to create links in the text
%%%%%%%%%%%%%%% Polish letter packages %%%%%%%%%%%%%%%%%%%%%%%%%%%%%%%%
\usepackage[polish]{babel}
\usepackage[utf8]{inputenc}
\usepackage{t1enc}
%%% package for generating gibberish (placeholder text):
\usepackage{lipsum}
%%% package for adding fictitious notes `a la ``todo''
\usepackage{todonotes}
% This package is supposedly intended, among other things, to allow
% line breaks in table cells and in the table ``header''
% (table environment)
\usepackage{makecell}
% this package is used for elegant formatting of listings (programs? but presumably not only programs)
\usepackage{listings}
% this package is used, among other things, to check whether the appropriate bibliography entries
% are actually cited in the body of the file. Presumably, it should
% be run only when checking the file; in the official version it should be disabled.
\usepackage{refcheck}
% package included for the time being so that we can create ``monstrous'' bigforall and bigexists quantifiers
\usepackage{graphicx}
% the rotating package is used for rotating objects such as text, tables, images, graphs, etc.
\usepackage{rotating}
% package for conveniently creating multi-row tables.
\usepackage{multirow}
% the first package makes it possible to use calligraphic fonts,
% whereas I have no idea what the purpose of the second package is,
% but it is certainly related to the first one.
\usepackage{calligra}
\usepackage[T1]{fontenc}

\newcommand{\countableEllentuckPerfect}{\mathrm{Count-Ell-Perf}}
\newcommand{\fcInItselfEllentuckPerfect}{\mathrm{FC-in-itself-Ell-Perf}}
\newcommand{\separableEllentuckPerfect}{\mathrm{Sep-Ell-Perf}}

\begin{document}
In the Ellentuck topology, we distinguish the following:
\begin{enumerate}
\item
A countable Ellentuck perfect set ($\countableEllentuckPerfect$)
\item
An Ellentuck perfect subset that is first category in itself
($\fcInItselfEllentuckPerfect$)
\item
An Ellentuck perfect separable (i.e. having a dense countable set) subset ($\separableEllentuckPerfect$).
\end{enumerate}

Obviously 
  $\countableEllentuckPerfect \implies \fcInItselfEllentuckPerfect$
and
  $\countableEllentuckPerfect \implies \separableEllentuckPerfect$
but there is no relationship between
$\fcInItselfEllentuckPerfect$ and $\separableEllentuckPerfect$

I do not know whether we explicitly wrote this in our paper on the dichotomy,
but the fact that the Ellentuck topology contains the Sorgenfrey line
as a perfect subset implies a counterexample to your
question from your first basic article.
I could not find it anywhere.

The two major results from your article:
\begin{itemize}
\item
  Every Ellentuck dense subset of $[\omega]^\omega$ contains
a set from $\countableEllentuckPerfect$;
\item
  Every Ellentuck dense subset of any $P \in \fcInItselfEllentuckPerfect$
contains a set from $\countableEllentuckPerfect$.
\end{itemize}

The question was:
Is the assumption that $P$ is first category in itself necessary?

The answer is YES, and an example can be constructed using the fact that the
Ellentuck space contains the Sorgenfrey line.

Thats because:

Wiemy że topologia Ellentucka zawiera 
Sorgenfrey line jako swój podzbiór doskonały.
O linii Sorgenfreya można dowieść taki lemat
(dowód jest bardzo łatwy):
{\bf Theorem: }\\
If $P \subseteq (0,1]$ is a Sorgenfrey perfect set, then
the euclidean closure $\mathit{cl}(P)$ is a classical perfect
set and $\mathit{cl}(P) \setminus P$ is a countable set.
\end{document}
